% !TeX root = ../../main.tex
\subsection{正比例について}

2つの数量 $x$，$y$ があり，$y$ が $x$ に比例するとき，

\[y = ax \qquad (a \neq 0)\]

という関係が成り立つ。このとき $a$ を\textbf{比例定数}という.

\begin{definitionbox}{\textbf{正比例}}
    2つの数量 $x$，$y$ があり，$y$ が $x$ に比例するとき，$y = ax$（$a \neq 0$）という関係が成り立つ.
    このとき $a$ を\textbf{比例定数}という.
\end{definitionbox}

\begin{theorembox}{\textbf{正比例のグラフ}}
    $y = ax$ のグラフは\textbf{原点を通る直線}である.
    \begin{enumerate}
        \item $a > 0$ のとき，グラフは\textbf{右上がり}である.
        \item $a < 0$ のとき，グラフは\textbf{右下がり}である.
        \item $A \leqq x \leqq B$ のとき，$y$ の変域は $C \leqq y \leqq D$ となる（$A \leqq x \leqq B$ の範囲で $y$ は $C$ から $D$ まで変化する）.
    \end{enumerate}
\end{theorembox}

\begin{example}[【例題 1】]
    $y$ が $x$ に比例し，$x = 3$ のとき $y = 6$ である. $y$ を $x$ の式で表しなさい.
\end{example}

\begin{proof}\mbox{}\\
    $y = ax$ とおく.

    \begin{hiddenanswer}
        $x = 3$，$y = 6$ を代入すると，
        \begin{align*}
            6 &= 3a \\
            a &= 2
        \end{align*}
        \finalanswer{答え：$y = 2x$}
    \end{hiddenanswer}
\end{proof}

\newpage
